\documentclass[otchet]{SCWorks}
% Тип обучения (одно из значений):
%    bachelor   - бакалавриат (по умолчанию)
%    spec       - специальность
%    master     - магистратура
% Форма обучения (одно из значений):
%    och        - очное (по умолчанию)
%    zaoch      - заочное
% Тип работы (одно из значений):
%    coursework - курсовая работа (по умолчанию)
%    referat    - реферат
%  * otchet     - универсальный отчет
%  * nirjournal - журнал НИР
%  * digital    - итоговая работа для цифровой кафедры
%    diploma    - дипломная работа
%    pract      - отчет о научно-исследовательской работе
%    autoref    - автореферат выпускной работы
%    assignment - задание на выпускную квалификационную работу
%    review     - отзыв руководителя
%    critique   - рецензия на выпускную работу

% * Добавлены вручную. За вопросами к @mchernigin
\usepackage{preamble}

\begin{document}

% Кафедра (в родительном падеже)
\chair{математической кибернетики и компьютерных наук}

% Тема работы
\title{Сравнение методов Д'Ондта и Сент-Лагю}

% Курс
\course{1}

% Группа
\group{151}

% Факультет (в родительном падеже) (по умолчанию "факультета КНиИТ")
% \department{факультета КНиИТ}

% Специальность/направление код - наименование
% \napravlenie{02.03.02 "--- Фундаментальная информатика и информационные технологии}
% \napravlenie{02.03.01 "--- Математическое обеспечение и администрирование информационных систем}
% \napravlenie{09.03.01 "--- Информатика и вычислительная техника}
\napravlenie{09.03.04 "--- Программная инженерия}
% \napravlenie{10.05.01 "--- Компьютерная безопасность}

% Для студентки. Для работы студента следующая команда не нужна.
% \studenttitle{студентки}

% Фамилия, имя, отчество в родительном падеже
\author{Захаркина Андрея Юрьевича}

% Заведующий кафедрой 
\chtitle{доцент, к.\,ф.-м.\,н.}
\chname{С.\,В.\,Миронов}

% Руководитель ДПП ПП для цифровой кафедры (перекрывает заведующего кафедры)
% \chpretitle{
%     заведующий кафедрой математических основ информатики и олимпиадного\\
%     программирования на базе МАОУ <<Ф"=Т лицей №1>>
% }
% \chtitle{г. Саратов, к.\,ф.-м.\,н., доцент}
% \chname{Кондратова\, Ю.\,Н.}

% Научный руководитель (для реферата преподаватель проверяющий работу)
\satitle{доцент, к.\,ф.-м.\,н.} %должность, степень, звание
\saname{Г.\,Г.\,Наркайтис}

% Руководитель практики от организации (руководитель для цифровой кафедры)
\patitle{доцент, к.\,ф.-м.\,н.}
\paname{С.\,В.\,Миронов}

% Руководитель НИР
\nirtitle{доцент, к.\,п.\,н.} % степень, звание
\nirname{В.\,А.\,Векслер}

% Семестр (только для практики, для остальных типов работ не используется)
\term{2}

% Наименование практики (только для практики, для остальных типов работ не
% используется)
\practtype{учебная}

% Продолжительность практики (количество недель) (только для практики, для
% остальных типов работ не используется)
\duration{2}

% Даты начала и окончания практики (только для практики, для остальных типов
% работ не используется)
\practStart{01.07.2024}
\practFinish{13.01.2024}

% Год выполнения отчета
\date{2026}

\maketitle

% Включение нумерации рисунков, формул и таблиц по разделам (по умолчанию -
% нумерация сквозная) (допускается оба вида нумерации)
\secNumbering

\tableofcontents

% Раздел "Обозначения и сокращения". Может отсутствовать в работе
\abbreviations
\begin{description}
     \item $m$ "--- кол-во партий (списков), участвующих в распределении;
     \item $V_i$ "---  число голосов, полученных партией $i$, где $i = 1, 2, \cdots,m$;
     \item $V = \sum_{i = 1}^{m}$ "--- общее число поданных голосов;
     \item $S$ "--- общее кол-во распределяемых мандатов;
     \item $s_i$ "---  количество мандатов, уже выделенных партии $i$ на текущем шаге алгоритма.
     
     Квота Хейра (определяет <<стоимость>> одного мандата): 
     \begin{equation*}
         Q = \frac{V}{S}
     \end{equation*}
\end{description}

% Раздел "Определения". Может отсутствовать в работе
% \definitions

% Раздел "Определения, обозначения и сокращения". Может отсутствовать в работе.
% Если присутствует, то заменяет собой разделы "Обозначения и сокращения" и
% "Определения"
% \defabbr

\intro
Пропорциональная избирательная система часто требует распределения ограниченного числа мест (мандатов) между несколькими партиями пропорционально поданным за них голосам. Поскольку число мест дискретно, а голоса измеряются целыми числами, точная пропорциональность почти всегда невозможна. Для решения этой задачи используются различные методы делителей, наиболее известные из которых — метод Д’Ондта и метод Сент-Лагю (в исходном варианте — метод Вебстера) \cite{balfair,puke}. Оба метода основаны на последовательном делении голосов партий на возрастающие числа, но используют разные последовательности делителей, что приводит к разной степени пропорциональности и по-разному влияет на крупные и мелкие партии. Кроме того, данные методы можно вполне использовать для округления любых величин с константной суммой \cite{gallprop}.

В статье приводится математическое описание обоих методов, три ключевые формулы, таблица с числовым примером и краткий анализ их политических следствий.
\section{Метод Д'Ондта}
Метод предложен бельгийским математиком Виктором Д’Ондтом в 1878 г. и известен также как <<метод делителей $1, 2, 3, \ldots$>>. Он обладает свойством благоприятствовать крупным партиям (завышать их представительство относительно точной пропорции) \cite{balfair,gallprop}.
\subsection{Алгоритм метода Д'Ондта}
\begin{enumerate}
    \item Для каждой партии $i$ вычисляется текущее значение
    \begin{equation*}
        D^{(k)}_i = \frac{V_i}{s_i+1}
    \end{equation*}
    \item Мандат получает та партия, у которой значение 
$ D^{(k)}_i$ максимально. При равенстве решает жребий или дополнительное правило.
\item У победившей партии $s_i$ увеличивается на 1, и алгоритм повторяется, пока все $S$ мандатов не будут распределены.
\end{enumerate}

Крупные партии получают мандаты чаще, чем предписывает строгая пропорция, потому что их голоса делятся на меньшие числа на ранних шагах. Тем временем мелкие партии с малым $V_i$ могут вообще не получить ни одного мандата, даже если их доля голосов превышала бы квоту $Q$, но не достигает выигрышных значений по сравнению с уже поделёнными голосами лидеров \cite{puke}.
\section{Метод Сент-Лагю}
Метод назван в честь французского математика Андре Сент-Лагю, хотя исторически эквивалентный метод был предложен сенатором Дэниэлом Вебстером для распределения мест в Конгрессе США \cite{balfair}. Он использует нечётные делители $1,3,5,\ldots$ и считается нейтральнее по отношению к размеру партий; в чистом виде он не даёт систематического преимущества ни крупным, ни мелким партиям \cite{gallprop}.

Алгоритм полностью аналогичен алгоритму метода Д'Ондта за исключением использования нечетных делителей:
\begin{equation*}
    D^{(k)}_i = \frac{V_i}{2s_i+1}
\end{equation*}

Например, если в выборах участовало 4 партии получившие соответственно 10000, 80000, 30000 и 20000 голосов, то распределение 8 мандатов для них будет выглядеть следующим образом:
\begin{center}
    \includegraphics[width = 0.7\textwidth]{LaTeX/task3/img/sl.png}

    Рисунок 1. На Википедии лежат эти красивые png-шки, ну я их и использую
\end{center}

\section{Сравнительный пример}
Теперь рассмотрим распределение 10 мандатов для 5 партий (A, B,C, D, E) c голосами 41000, 29000, 15000, 10000 и 5000 соответственно обоими методами.
\begin{table}[h]
\caption{Сравнительный пример методов Сент-Лагю и Д'Ондта}
\label{tab:comparison}
\begin{tabular}{|c|c|c|c|c|}
\hline
Партия & Голоса, \% & Точная пропорция & Д'Ондт & Сент-Лагю \\
\hline
A & 41.0 & 4.1 & 5 & 4 \\
\hline
B & 29.0 & 2.9 & 3 & 3 \\
\hline
C & 15.0 & 1.5 & 1 & 1 \\
\hline
D & 10.0 & 1 & 1 & 1 \\
\hline
E & 5.0  & 0.5 & 0 & 1 \\
\hline
Итого & 100.0 & 10 & 10 & 10 \\
\hline
\end{tabular}
\end{table}

\begin{center}
    
\includegraphics[width = 0.7\textwidth]{LaTeX/task3/img/Screenshot 2026-05-22 024410.png}}

    Рисунок 2. Распределение методом Д'Ондта. Шакально конечно, но картинки на Википедии нет, пришлось делать самому, а .svg латех не поддерживает
\end{center}

\begin{center}
    \includegraphics[width = 0.7\textwidth]{LaTeX/task3/img/Screenshot 2026-05-22 025556.png}

    Рисунок 3. Распределение методом Сент-Лагю
\end{center}
Как видно из таблицы 3.1 метод Д'Ондта завысил результаты лидирующей партии, а партии E вообще не досталось мест, в то время как метод Сент-Лагю более близок к точному кол-ву <<заслуженных>> мест.
\section{Код}
Собственно sainte\_lague.cpp и является причиной, почему я выбрал именно эту тему для задания. Я положил его в репозиторий по приколу, чтобы был, и так он и лежал в корне один, никуда не принадлежащий, ни рыба ни мясо. Вот я и решил найти ему применение и перенести в какую-нибудь директорию.

Что касается самого кода, выглядит он довольно ужасно. Как понятно из названия, изначально программа предполагалась для расчета только метода Сент-Лагю, но потом я решил, что для сравнения лучше сразу выводить два метода рядом. Отсюда и дублированние переменные по типу \texttt{seat\_num1},\texttt{extra\_votes1} и т.д.

\inputminted[fontsize=\small, breaklines=true, style=borland, linenos]{cpp}{sainte_lague.cpp}
% После введения — серии \section, \subsection и т.д.

\conclusion
Методы делителей Д’Ондта и Сент-Лагю решают одну и ту же задачу дискретного пропорционального распределения, но дают разный политический эффект. Первый, используя делители  $1,2,3,\ldots$ благоприятствует крупным партиям; второй, с делителями 
$1,3,5,\ldots$ обеспечивает почти идеальную пропорциональность без систематического смещения. Выбор между ними — это всегда компромисс между стабильностью правительства и точностью представительства \cite{balfair}. Понимание математических основ позволяет осознанно подходить к распределению мандатов и не только.
% Отобразить все источники. Даже те, на которые нет ссылок.
 \nocite{*}

\bibliographystyle{ugost2003}
\bibliography{thesis}

% Окончание основного документа и начало приложений Каждая последующая секция
% документа будет являться приложением
\appendix

\end{document}
